\vspace{0.5em}\noindent\textbf{AWS Networking Advanced: 3 Techniques for Building an
Enterprise-Grade
VPC}\par

\vspace{0.5em}\noindent\textbf{Introduction}\par

A basic Amazon VPC architecture consisting of public subnets, private
subnets, and a NAT Gateway is a suitable foundation for deploying
applications on AWS. This architecture separates resources that require
direct Internet access from internal components such as application
servers, databases, and caches.

However, when a system expands to include Development, Staging, and
Production environments or is divided across multiple microservices and
VPCs, the basic architecture may begin to present three major
limitations:

\begin{itemize}
\tightlist
\item
  Increasing NAT Gateway costs caused by frequent access from private
  subnets to Amazon S3, Amazon DynamoDB, and other AWS services.
\item
  A potential single point of failure when all outbound traffic depends
  on one NAT Gateway in one Availability Zone.
\item
  Complex network management when many VPCs are connected through a
  full-mesh VPC Peering model.
\end{itemize}

This article presents three techniques for upgrading a VPC architecture
for enterprise environments: VPC Endpoints, Multi-AZ networking, and AWS
Transit Gateway.

\vspace{0.5em}\noindent\textbf{1. Reducing Costs and Improving Security with VPC
Endpoints}\par

\vspace{0.5em}\noindent\textbf{The NAT Gateway Traffic
Problem}\par

EC2 instances, containers, and application servers running in private
subnets usually do not have public IP addresses. When these resources
need to access Amazon S3, DynamoDB, or other AWS services, their traffic
may pass through a NAT Gateway.

A basic traffic flow may look like this:

\begin{Shaded}
\begin{Highlighting}[]
\NormalTok{Private Subnet}
\NormalTok{      |}
\NormalTok{      v}
\NormalTok{NAT Gateway}
\NormalTok{      |}
\NormalTok{      v}
\NormalTok{AWS Public Service Endpoint}
\end{Highlighting}
\end{Shaded}

This design may introduce hourly NAT Gateway charges and data-processing
costs. Traffic may also pass through a public service endpoint even
though both the application and destination service operate within AWS
infrastructure.

\vspace{0.5em}\noindent\textbf{Gateway Endpoints}\par

Gateway Endpoints are available for two primary services:

\begin{itemize}
\tightlist
\item
  Amazon S3.
\item
  Amazon DynamoDB.
\end{itemize}

After a Gateway Endpoint is created and associated with the route tables
of private subnets, traffic to S3 or DynamoDB can travel directly
through the AWS network without using a NAT Gateway.

The updated traffic flow can be represented as follows:

\begin{Shaded}
\begin{Highlighting}[]
\NormalTok{Private Subnet}
\NormalTok{      |}
\NormalTok{      v}
\NormalTok{VPC Gateway Endpoint}
\NormalTok{      |}
\NormalTok{      v}
\NormalTok{Amazon S3 or Amazon DynamoDB}
\end{Highlighting}
\end{Shaded}

Gateway Endpoints provide several benefits:

\begin{itemize}
\tightlist
\item
  A NAT Gateway is no longer required for S3 and DynamoDB traffic.
\item
  Traffic remains within the AWS network.
\item
  NAT Gateway data-processing costs can be reduced.
\item
  Access can be controlled through endpoint policies.
\item
  Workloads become less dependent on outbound Internet connectivity.
\end{itemize}

Gateway Endpoints for S3 and DynamoDB do not have hourly endpoint
charges, which makes them one of the first options to consider when
optimizing network costs.

\vspace{0.5em}\noindent\textbf{Interface Endpoints and AWS
PrivateLink}\par

For AWS services that do not support Gateway Endpoints, an Interface
Endpoint can be used.

An Interface Endpoint creates Elastic Network Interfaces inside selected
subnets and assigns private IP addresses to them. Applications can then
access the target AWS service through private IP addresses or private
DNS without using the public Internet.

Services commonly accessed through Interface Endpoints include:

\begin{itemize}
\tightlist
\item
  Amazon CloudWatch.
\item
  AWS Secrets Manager.
\item
  Amazon ECR.
\item
  AWS Systems Manager.
\item
  AWS Key Management Service.
\item
  Amazon Kinesis.
\end{itemize}

Interface Endpoints use AWS PrivateLink to keep service traffic inside
the AWS network.

Compared with NAT Gateway access, this design can:

\begin{itemize}
\tightlist
\item
  Reduce outbound Internet traffic.
\item
  Control access through Security Groups.
\item
  Reduce the attack surface.
\item
  Support a clearer private-networking model.
\item
  Allow private workloads to access AWS services through private IP
  addresses.
\end{itemize}

Interface Endpoints still have hourly and data-processing charges. Their
cost should therefore be compared with NAT Gateway costs based on the
actual traffic patterns of the system.

\vspace{0.5em}\noindent\textbf{2. Deploying Across Multiple Availability Zones and Avoiding
Unnecessary Cross-AZ
Costs}\par

\vspace{0.5em}\noindent\textbf{Why Multiple Availability Zones
Matter}\par

An Availability Zone can experience an independent failure. If all
application servers, databases, and outbound networking resources are
placed in one AZ, the entire system may become unavailable when that AZ
is affected.

For higher availability, production architectures are commonly deployed
across at least two Availability Zones.

For example:

\begin{Shaded}
\begin{Highlighting}[]
\NormalTok{Availability Zone A}
\NormalTok{{-} Public Subnet A}
\NormalTok{{-} Private Subnet A}
\NormalTok{{-} NAT Gateway A}
\NormalTok{{-} Application instances A}

\NormalTok{Availability Zone B}
\NormalTok{{-} Public Subnet B}
\NormalTok{{-} Private Subnet B}
\NormalTok{{-} NAT Gateway B}
\NormalTok{{-} Application instances B}
\end{Highlighting}
\end{Shaded}

In this model, each Availability Zone has an independent outbound path.
If one AZ becomes unavailable, workloads in the remaining AZ can
continue using the NAT Gateway and resources available in that AZ.

\vspace{0.5em}\noindent\textbf{Avoiding a Single Point of
Failure}\par

If the system uses only one NAT Gateway in AZ-A, workloads in AZ-B must
send their outbound traffic across Availability Zone boundaries before
reaching the gateway.

This creates two issues:

\begin{enumerate}
\def\labelenumi{\arabic{enumi}.}
\tightlist
\item
  The NAT Gateway in AZ-A becomes a shared point of dependency.
\item
  Traffic from AZ-B to AZ-A may generate Cross-AZ data-transfer charges.
\end{enumerate}

A more resilient architecture deploys one NAT Gateway in every active
Availability Zone. Each private subnet route table is configured to use
the NAT Gateway located in the same AZ.

For example:

\begin{Shaded}
\begin{Highlighting}[]
\NormalTok{Private Subnet A {-}\textgreater{} NAT Gateway A}
\NormalTok{Private Subnet B {-}\textgreater{} NAT Gateway B}
\end{Highlighting}
\end{Shaded}

If AZ-A becomes unavailable, resources in AZ-B still have an independent
outbound path and do not depend on NAT Gateway A.

\vspace{0.5em}\noindent\textbf{AZ-Aware Routing}\par

AZ-aware routing prioritizes communication between workloads and
dependencies located in the same Availability Zone.

Examples include:

\begin{itemize}
\tightlist
\item
  An application instance in AZ-A using NAT Gateway A.
\item
  A workload in AZ-B connecting to a cache node in AZ-B.
\item
  A load balancer distributing requests to suitable targets.
\item
  Database replicas or cache nodes being distributed across multiple
  AZs.
\item
  Service discovery preferring endpoints located closer to the workload
  when supported.
\end{itemize}

This design can:

\begin{itemize}
\tightlist
\item
  Reduce unnecessary Cross-AZ traffic.
\item
  Reduce network latency.
\item
  Limit inter-AZ data-transfer costs.
\item
  Allow each AZ to operate more independently.
\end{itemize}

However, AZ locality must still be balanced with High Availability
requirements. An application should not depend entirely on one local
resource if that dependency does not provide a failover mechanism.

\vspace{0.5em}\noindent\textbf{3. Managing Multi-VPC Connectivity with AWS Transit
Gateway}\par

\vspace{0.5em}\noindent\textbf{Limitations of VPC
Peering}\par

VPC Peering works well when only a small number of VPCs need to
communicate. Two VPCs can be connected directly and exchange traffic
through private IP addresses.

However, VPC Peering does not support transitive routing.

For example:

\begin{Shaded}
\begin{Highlighting}[]
\NormalTok{VPC{-}A \textless{}{-}\textgreater{} VPC{-}B}
\NormalTok{VPC{-}B \textless{}{-}\textgreater{} VPC{-}C}
\end{Highlighting}
\end{Shaded}

The fact that VPC-A connects to VPC-B and VPC-B connects to VPC-C does
not automatically allow VPC-A to communicate with VPC-C.

When many VPCs require direct connections, the number of VPC Peering
connections can grow according to the following formula:

\begin{Shaded}
\begin{Highlighting}[]
\NormalTok{N(N {-} 1) / 2}
\end{Highlighting}
\end{Shaded}

With 10 VPCs, a full-mesh design may require up to 45 peering
connections.

At this scale, managing route tables, CIDR ranges, Security Groups, and
network policies becomes increasingly difficult. Adding one new VPC may
also require multiple new connections and route-table updates.

\vspace{0.5em}\noindent\textbf{Transit Gateway as a Cloud
Router}\par

AWS Transit Gateway acts as a centralized router for multiple VPCs and
external networks.

Resources that can connect to a Transit Gateway include:

\begin{itemize}
\tightlist
\item
  Amazon VPCs.
\item
  Site-to-Site VPN connections.
\item
  AWS Direct Connect through a Direct Connect Gateway.
\item
  VPCs from multiple AWS accounts.
\item
  Development, Staging, and Production environments.
\item
  Shared Services VPCs.
\item
  On-Premises data center networks.
\end{itemize}

Instead of creating direct connections between every pair of VPCs, each
VPC creates a single attachment to the Transit Gateway.

\begin{Shaded}
\begin{Highlighting}[]
\NormalTok{             Development VPC}
\NormalTok{                    |}
\NormalTok{                    |}
\NormalTok{Shared VPC {-}{-}{-}{-} Transit Gateway {-}{-}{-}{-} Production VPC}
\NormalTok{                    |}
\NormalTok{                    |}
\NormalTok{             On{-}Premises Network}
\end{Highlighting}
\end{Shaded}

This hub-and-spoke model simplifies the network topology and centralizes
routing management.

\vspace{0.5em}\noindent\textbf{Network Segmentation with Transit Gateway Route
Tables}\par

Transit Gateway can use multiple route tables to control which VPCs and
networks are allowed to communicate.

For example:

\begin{itemize}
\tightlist
\item
  The Development VPC can access the Shared Services VPC.
\item
  The Production VPC can access the Shared Services VPC.
\item
  The Development VPC cannot directly access the Production VPC.
\item
  The On-Premises Network can access only selected subnets.
\item
  A Security VPC can inspect traffic before forwarding it to other VPCs.
\end{itemize}

This segmentation can:

\begin{itemize}
\tightlist
\item
  Reduce lateral movement between environments.
\item
  Separate Development, Staging, and Production more clearly.
\item
  Centralize network-policy management.
\item
  Simplify the addition of new VPCs.
\item
  Support multi-account architectures with AWS Organizations.
\end{itemize}

\vspace{0.5em}\noindent\textbf{Comparing Basic and Enterprise VPC
Architectures}\par

{
\begingroup
\small
\setlength{\tabcolsep}{3pt}
\renewcommand{\arraystretch}{1.15}
\begin{longtable}{@{}
  >{\raggedright\arraybackslash}p{(\linewidth - 4\tabcolsep) * \real{0.3333}}
  >{\raggedright\arraybackslash}p{(\linewidth - 4\tabcolsep) * \real{0.3333}}
  >{\raggedright\arraybackslash}p{(\linewidth - 4\tabcolsep) * \real{0.3333}}@{}}
\toprule\noalign{}
\begin{minipage}[b]{\linewidth}\raggedright
Category
\end{minipage} & \begin{minipage}[b]{\linewidth}\raggedright
Basic Architecture
\end{minipage} & \begin{minipage}[b]{\linewidth}\raggedright
Enterprise Architecture
\end{minipage} \\
\midrule\noalign{}
\endhead
\bottomrule\noalign{}
\endlastfoot
Access to S3 and DynamoDB & Through a NAT Gateway or public endpoint &
Through a Gateway Endpoint \\
Access to other AWS services & Commonly through a NAT Gateway & Through
Interface Endpoints and AWS PrivateLink when appropriate \\
High Availability & May depend on one NAT Gateway & Independent NAT
Gateway in each Availability Zone \\
Cross-AZ traffic & May increase because of inefficient routing & Reduced
through AZ-aware routing \\
Multi-VPC connectivity & VPC Peering configured manually between pairs &
Centrally managed through Transit Gateway \\
Development and Production separation & Depends on multiple routes and
peering connections & Controlled through Transit Gateway route tables \\
Scalability & Suitable for smaller systems & More suitable for multi-VPC
and multi-account systems \\
Management & Simple initially but difficult to control at scale & More
complex initially but easier to manage centrally \\
\end{longtable}
\endgroup
}

\vspace{0.5em}\noindent\textbf{When Should Each Technique Be
Used?}\par

\vspace{0.5em}\noindent\textbf{Use VPC Endpoints when:}\par

\begin{itemize}
\tightlist
\item
  Workloads in private subnets frequently access Amazon S3 or DynamoDB.
\item
  Applications need to access AWS services without using the public
  Internet.
\item
  NAT Gateway traffic and processing costs need to be reduced.
\item
  The system has stronger data-security requirements.
\item
  Access needs to be limited through endpoint policies.
\end{itemize}

\vspace{0.5em}\noindent\textbf{Use Multi-AZ NAT Gateways
when:}\par

\begin{itemize}
\tightlist
\item
  The system requires High Availability.
\item
  Workloads operate across multiple Availability Zones.
\item
  A single NAT Gateway should not become a single point of failure.
\item
  Unnecessary Cross-AZ traffic should be reduced.
\item
  Outbound connectivity must remain available when one AZ fails.
\end{itemize}

\vspace{0.5em}\noindent\textbf{Use AWS Transit Gateway
when:}\par

\begin{itemize}
\tightlist
\item
  The organization operates many VPCs or AWS accounts.
\item
  On-Premises networks must connect to AWS.
\item
  VPC Peering has become difficult to manage.
\item
  Development, Staging, and Production require clear separation.
\item
  Routing needs to be controlled from a central location.
\item
  Shared Services or centralized Security Inspection architectures are
  required.
\end{itemize}

\vspace{0.5em}\noindent\textbf{Cost Considerations}\par

An advanced network architecture does not mean that every component
becomes less expensive.

Gateway Endpoints for S3 and DynamoDB can reduce traffic through NAT
Gateways. However, Interface Endpoints, NAT Gateways, and Transit
Gateway each have their own pricing models.

Before implementation, the architecture should consider:

\begin{itemize}
\tightlist
\item
  Total monthly data-transfer volume.
\item
  AWS services frequently accessed by private workloads.
\item
  Number of Availability Zones.
\item
  Number of VPCs and Transit Gateway attachments.
\item
  Frequency of Cross-AZ communication.
\item
  Internet outbound traffic.
\item
  Security and availability requirements.
\item
  Operational cost and management complexity.
\end{itemize}

The objective is not simply to minimize cost. It is to find an
appropriate balance between cost, security, performance, scalability,
and High Availability.

\vspace{0.5em}\noindent\textbf{Conclusion}\par

The combination of public subnets, private subnets, and a NAT Gateway
provides a strong starting point for building applications on AWS. As
the system grows across multiple environments, VPCs, or AWS accounts,
the network architecture must also evolve.

Three important techniques are:

\begin{enumerate}
\def\labelenumi{\arabic{enumi}.}
\tightlist
\item
  Use VPC Endpoints to keep AWS service traffic inside the AWS network
  and reduce dependence on NAT Gateways.
\item
  Deploy independent NAT Gateways across multiple Availability Zones to
  improve availability and reduce unnecessary Cross-AZ traffic.
\item
  Use AWS Transit Gateway to manage connectivity between multiple VPCs
  through a centralized model.
\end{enumerate}

VPC optimization is not only a routing problem. It also involves
balancing security, scalability, availability, performance, and cost.

A strong Enterprise VPC Architecture should be based on the actual
requirements of the system. Additional services should not be introduced
only to make the architecture appear more complex without providing
clear value.

\vspace{0.5em}\noindent\textbf{Publication Information}\par

\begin{itemize}
\tightlist
\item
  \textbf{Topic:} Upgrading Amazon VPC architecture for enterprise
  systems.
\item
  \textbf{Published date:} July 29, 2026.
\item
  \textbf{Platform:} AWS Study Groups.
\item
  \textbf{Status:} Published.
\item
  \textbf{Public link:}
  \href{https://www.facebook.com/groups/awsstudygroupfcj/permalink/2225997548165205/\#}{Blog
  4 on AWS Study Groups}
\end{itemize}
